\section{introduction}\label{sec:introduction}

\subsection{Motivation}\label{subsec:motivation}

Search and Rescue (SAR) operations are characterized by their inherent urgency and uncertainty.
Empirical analyses of SAR missions indicate that the likelihood of survival declines sharply as search time increases, to the point that missions may eventually become futile if continued indefinitely.
In the context of maritime incidents, survival and resuscitation probabilities are strongly linked to the duration of time spent submerged or exposed.
For example, ~\cite{RN2}, states that the probability of survival/resuscitation decreases significantly after 30 minutes of submersion at >6°C, and after 90 minutes at <6°C\@.

Traditional Maritime Search and Rescue (MSAR) relies on crewed vehicles like helicopters, seaborne vessels, and aircraft.
Keeping these assets at ``full readiness'', which in the context of MSAR means they must be adequately maintained, crewed, fuelled, and immediately deployable, can be difficult if not economically impossible in some regions.
Delays due to limited readiness directly translate into an increase in response time, as time that should be spent on the operation is spent remediating issues at sea, which could lead to a decrease in survival rates.

Research into the Human operators indicates the cognitive and physical demands of MSAR operations can lead to degraded performance.
\cite{riverine_MSAR_stress} directly reports that conditions combining higher physiological stress (like an increased heart rate) and distraction lead to lowered search success.
The paper interprets this as MSAR performance, at least in a riverine context, is significantly impacted by distraction, fatigue, and overwork.
Operators reporting neither stress nor exertion located 78\% of deployed targets, while crew reporting both conditions located only 57\%.
Operators reporting only Exertion located 70\% of targets, while those reporting only Distraction located 57\% of targets.
While the paper stops short of claiming a perfectly isolated causal relationship due to factors like experience, and other pre-SAR factors, it does suggest that the physical and cognitive demands of MSAR are impacted by the conditions of the operators, and that these conditions could be improved leading to better performance from the operators.

\cite{7761074} shows that Unmanned Aerial Vehicles(UAVs) and Unmanned Surface Vessels(USVs) can be used to augment MSAR operations, by providing rapid aerial and surface coverage.
However, this system remained non-autonomous or semi-autonomous, relying on a team of human operators to control the drones either one-to-one or in small groups.
This would occur due to regulatory constraints, safety concerns, and incomplete trust in onboard autonomy.

As the number of involved drones increases, these control schemes would struggle to scale, facing manpower issues, and cognitive workload from inter-operator coordination.
Digital Twin (DT) frameworks promise to ease these issues by aggregating state information, predicting behaviour, and supporting decision-making within a single interactive model.

BEACON (Behavioral Engine for Autonomous Coordination in Ocean Navigation) aims to explore how DT-inspired architectures and HCI principles can be combined to support a single operator supervising a co-ordinated swarm of autonomous drones, in an MSAR-like scenario.

\subsection{Problem Statement}\label{subsec:problem-statement}
The central problem this project seeks to address is the challenge of effectively supervising an autonomous drone swarm in MSAR-like scenarios, while not incurring unacceptable cognitive workloads or losing situation awareness.

This central problem can be broken down into the following sub-problems:

\begin{itemize}
    \item How can we design a simulation and C3I (Command, Control, Communications, and Intelligence) framework captures key MSAR constraints in a way that can evolve into a complete Digital Twin system?
    \item How can we design an interface that reveals the structure of the search problem and exposes high-level swarm autonomy decisions while avoiding information overload?
    \item What metrics and evaluation methods do we use to evaluate the effectiveness of the system?
\end{itemize}

\subsection{Contributions}\label{subsec:contributions}

This project seeks to contribute, through the design and implementation of BEACON, a five-layer Digital Twin framework for MSAR drone swarms.
This would be implemented as a React/TypeScript simulation.

That simulation needs to include a workload-aware supervisory interface, with a top-down mission map, a four-tier EICAS-style alert hierarchy, and HCI-grounded design principles.
This would be followed up with an evaluation framework, with explicit rationale for metrics selection and exclusions.
Finally, this would include an analysis of prior work on MSAR, Drone Swarms, HCI in Supervisory systems, and workload measurement, to challenge common assumptions and identify BEACON's scope and limitations.

While not intended to directly replace current operational MSAR systems, BEACON seeks to inform their design.
The expected impact is to ground BEACON's interface in Situational Awareness, workload theory, and HCI principles, and empirically evaluate the system's effectiveness in reducing operator workload and improving situational awareness, compared to traditional methods.

BEACON's clear separation between the domain logic and its UI, in addition to its deterministic environment generation, and structured metrics collection, should make it a suitable testbed for future HCI implementations in Human-Drone Swarm Interfaces, Automation Transparency, and stress management.

If integrated into a larger-scale research programme, BEACON could enable systematic, controlled testing of human-swarm MSAR interfaces before being deployed in costly real-world scenarios.
